% !TEX root = chapter-standalone.tex

\section{Generators and relations}
\label{sec:generators-relations-algebra}

In this section we develop an economical way of describing \SY{semigroups} and \SY{monoids}: instead of listing all elements and tabulating the whole composition operation, we specify a few \emph{generators}, from which all elements can be built by composition, together with some \emph{relations}, which record the equations that hold between such composites.
We first examine these two ideas in structures we already have; then, reversing the perspective, we show how to \emph{define} new \SY{semigroups} and \SY{monoids} by prescribing generators and relations from scratch---a construction that we will put to work in \cref{sec:sums-algebra}.

\subsection{Generating subsets}

In \cref{exa:plant-trafo-semigroup} we considered a set of states
%
\begin{equation}
    \setA= \makeset{ \sprout, \yng, \mature, \old, \dead },
\end{equation}
%
a function~$\mapa \colon \setA \sto \setA$, and the semigroup
%
\begin{equation}
    \label{eq:trafo-sgrp-one-gen}
    \sgrpA = \makeset{ \mapa^n \mid n \setin \natnumbers }.
\end{equation}
%
Note that~$\sgrpA$ has a special form: all of its elements can be expressed in terms of one of its elements,~$\mapa$, and the multiplication operation (which, in this case, is function composition).
To describe this state of affairs we say that~$\sgrpA$ is \emph{generated} by the element~$\mapa$.

\begin{ctdefinition}[Generating subsets]
    \label{def:gen-semigrp}
    Let~$\sgrpA = \tup{\sgrpAset, \mtimes}$ be a \SY{semigroup}, and let~$\setA \setsubseteq \sgrpAset$ be a subset.
    We say that~$\sgrpA$ is \emph{generated} by~\setA if every element of~$\sgrpA$ can be expressed as a finite composition of elements of~\setA.
\end{ctdefinition}

\begin{remark}
    Mutatis mutandis, the same definition also holds for \SY{monoids}.
    For \SY{groups}, we say \setA generates the group if every element of the group can be expressed as a finite composition of elements of \setA or their inverses.
\end{remark}

\begin{example}
    Consider \cref{exa:string-semigroup}, where elements of the \SY{semigroup}~$\sgrpA$ were non-empty lists built using the elements of the ``alphabet'' set~$\setA = \makeset{ \alphabeta, \alphabetb}$.
    In this case,~$\sgrpA$ is generated by~\setA.
\end{example}

\begin{example}
    Consider the natural numbers (without zero) as a \SY{semigroup}, where addition is the \SY{semigroup} composition operation (see \cref{exa:natnum-semigroup}).
    This \SY{semigroup} is generated by the subset~$\makeset{1 }$.
\end{example}

\subsection{Relations}

Let us return to the \SY{semigroup} \cref{eq:trafo-sgrp-one-gen}.
Recall that~$\mapa$ was defined by
%
\begin{align}
    \mapa(\sprout) & = \yng, \\
    \mapa(\yng)    & = \mature, \\
    \mapa(\mature) & = \old, \\
    \mapa( \old)   & = \dead, \\
    \mapa(\dead)   & = \dead.
\end{align}

Observe that the function~$\mapa^4$ will map all elements of~\setA to the element ``\dead''.
For example, if we start with the element ``\sprout'', the result of applying~$\mapa$ four times is
%
\begin{equation}
    \sprout \overset{\mapa}{\sto} \yng \overset{\mapa}{\sto} \mature \overset{\mapa}{\sto} \old \overset{\mapa}{\sto} \dead.
\end{equation}

Note also that for \emph{any}~$n \geq 4$, the function~$\mapa^n$ will map all elements of~\setA to the element ``\dead''.
If we consider~$\mapa^6$, for example, then, for any~$\ela \setin \setA$,
%
\begin{equation}
    \mapa^6(\ela) = \mapa^2(\mapa^4(\ela)) = \mapa^2(\dead) = \mapa(\mapa(\dead)) = \mapa(\dead) = \dead.
\end{equation}
%
It follows that all~$\mapa^n$, for~$n \geq 4$, are actually \emph{all the same map}: the one that sends every state to the dead state.
Thus, ~$\sgrpA = \makeset{ \mapa^n \mid n \setin \natnumbers }$ actually only has at most \emph{four} elements: $\mapa$,~$\mapa^2$,~$\mapa^3$, and~$\mapa^4$.

\begin{gradedexercise}[\exname{CheckRelations}]
    \label{ex:CheckRelations}
    Are any of the four maps~$\mapa$,~$\mapa^2$,~$\mapa^3$, and~$\mapa^4$ actually equal?
    Justify your answer by argumentation or by explicitly checking via calculation.
\end{gradedexercise}
\solutionof{CheckRelations}

When two elements which a priori could be distinct from each other (such as~$\mapa^6$ and~$\mapa^4$ above, for example) turn out to be equal, we call this a \emph{relation} between the elements of~$\sgrpA$.

\begin{ctdefinition}\label{def:relation-on-semigroup}
    \SYNDEF{relation on a semigroup}
    A \emph{relation} on a \SY{semigroup}~$\tup{\sgrpA, \mtimes}$ is an equation between compositions of elements of~$\sgrpA$.
\end{ctdefinition}

\begin{remark}
    Again, we have analogous definitions for \SY{monoids} and \SY{groups}.
    In these cases, we interpret the \SY{neutral element} $\idgrp$ as a ``zero-fold'' multiplication, so it can also be part of equations that express relations.
\end{remark}

\begin{remark}
    A word of warning on terminology: a ``relation on a \SY{semigroup}'' in the sense of \cref{def:relation-on-semigroup} is not the same thing as a binary \SY{relation} in the general sense of \cref{def:binary-relation}.
    The two notions are connected, however.
    An equation between composites of elements can be recorded as a \emph{pair} of composites (the left-hand side and the right-hand side), and a collection of such pairs is a binary endorelation on the underlying set of the \SY{semigroup}.
    Asking for the equations to actually \emph{hold} then amounts to asking that each such pair be contained in a suitable equivalence relation---in fact, in a \SY{congruence}, in the sense of \cref{def:congruence-semigroup}.
    We will exploit this point of view in \cref{sub:presentations} below.
\end{remark}

\begin{example}
    For the \SY{semigroup} \cref{eq:trafo-sgrp-one-gen}, the relations $\mapa^5 = \mapa^4$, $\mapa^6 = \mapa^5$, and $\mapa^6 = \mapa^4$, \etc are satisfied.
    However, the relation $\mapa^3 = \mapa$ is not satisfied, for example.

\end{example}

\begin{example}
    Consider the \SY{semigroup}~$\tup{\natnumbers, +}$.
    The equation~$l + k = k + l$ is an example of a relation that holds for all $l, k \setin \natnumbers$.
\end{example}

\begin{example}
    Consider the group $\grpA$ discussed in \cref{grp-order-three}, where
    %
    \begin{equation}
        \grpAset = \makesetBig{ 1, e^{\frac{1}{3}2\pi i}, e^{\frac{2 }{3}2\pi i} } \setsubseteq \cnumbers
    \end{equation}
    and the composition operation is multiplication of complex numbers.
    The element $\ela \definedas e^{\frac{1}{3}2\pi i}$ satisfies the relation $\ela^3 = \idgrp_\grpA$.
\end{example}

\begin{example}
    Consider the group $\grpA$ given in \cref{exa:grp-Klein4} which describes symmetries of a rectangle.
    We had
    \begin{equation}
        \grpAset = \makeset{ I, V, H, R}
    \end{equation}
    where $I = \idgrp$ corresponds to ``doing nothing'', $V$ is reflecting the rectangle along its long axis, $H$ is reflecting on the short axis, and $R$ is rotation by 180 degrees.

    In this group, the relations $V^2 = I$, $H^2 = I$, $R^2 = I$ are satisfied, for example.
\end{example}

\subsection{Freeness}

When, in \cref{def:gen-semigrp} we spoke about the \SY{semigroup}~$\sgrpA = \tup{\sgrpAset, \mtimes}$ being generated by a subset~$\setA \setsubseteq \sgrpAset$, we supposed that we already had a \SY{semigroup}~$\sgrpA$ to work with.
However, if we start with just a set, say~$\setA = \makeset{ \alphabeta, \alphabetb}$, then we saw in \cref{string-sgrp} that we can build a \SY{semigroup} from this set by considering lists of elements of~\setA, with concatenation as the composition operation.
The resulting \SY{semigroup} in that example has a special characteristic: its elements do not satisfy any relations other than the ones that are required by the definition of a \SY{semigroup}, namely those relations dictated by the \SY{associative law}.
Such a \SY{semigroup} is called \emph{free}.
If we think of relations as ``constraints'' (they are equations) between the elements of a \SY{semigroup}, then \SY{free semigroups} are ``free of constraints''.

For a given set, say~$\setA = \makeset{ \alphabeta, \alphabetb}$, there will in general be different ways of formally building a \SY{free semigroup} from it.
For instance, instead of considering lists of elements of~\setA
\begin{equation}
    \makelist{\alphabeta,\alphabeta, \alphabetb, \alphabeta, \alphabetb, \alphabetb, \alphabetb, \alphabeta},
\end{equation}
we could instead consider strings of elements
\begin{equation}
    \alphabeta\alphabeta \alphabetb \alphabeta \alphabetb \alphabetb \alphabetb \alphabeta,
\end{equation}
or tuples of elements
\begin{equation}
    \tup{\alphabeta, \alphabeta, \alphabetb, \alphabeta, \alphabetb, \alphabetb, \alphabetb, \alphabeta},
\end{equation}
both of which could also be composed in a way which is analogous to concatenation.

A common feature of all three of these formalizations is that we are writing a finite sequence of elements of~\setA, keeping account of the ordering.
Both approaches will in fact build a \SY{semigroup} from the set~\setA which is \emph{free}.
And in both cases there is a natural way of seeing~\setA as \emph{generating} the resulting \SY{semigroup}.
However, the two set-ups are \emph{formally} distinct because we are using a different way of writing things down with symbols.
We will see in a later chapter that the resulting \SY{semigroups} are essentially ``the same'' (they are isomorphic) and their ``freeness'' can be given an elegant characterization.

Because in fact \emph{all} \SY{free semigroups} constructed from~\setA are ``the same'', independent of the formal symbolic specifics of how they are constructed, we refer to them all as \emph{the} free \SY{semigroup} generated by~\setA.
Furthermore, if we were to work with a set~$\setB = \makeset{ \setAel, \setBel }$ instead of~$\setA = \makeset{ \alphabeta, \alphabetb}$, and generate a free \SY{semigroup} from~\setB, then this would also produce a \SY{semigroup} which is ``the same''.
Therefore, sometimes one speaks simply of ``the free semigroup on two generators'', in view of the fact that the essential feature is that both~\setA and~\setB have two elements.

Note that although~\setA and~\setB each only have two elements, the \SY{free semigroups} that they generate will have infinitely many elements.
Indeed, there are infinitely many lists that we can build from the elements of~\setA.
As we concatenate lists, the resulting compositions grow longer and longer, and there are no relations which would allow us to ``simplify'' a string to one which is shorter.

\begin{remark}
    The same discussion applies to \SY{monoids}: the set~$\listsof{\setA}$ of \emph{all} lists (including the empty list) with concatenation is the \emph{free monoid} on~$\setA$, as anticipated in \cref{exa:string-monoid}.
    Its elements satisfy no relations beyond those imposed by the \SY{monoid} axioms: associativity, and the neutrality laws for the empty list.
\end{remark}

\subsection{Defining structures by generators and relations}
\label{sub:presentations}

So far, generators and relations were features that we \emph{observed} in \SY{semigroups} and \SY{monoids} that we already had.
We now reverse the perspective.
Suppose we are handed only the raw data:
\begin{enumerate}
    \item a set~$\setA$ of symbols, which we intend as names for the generators, and
    \item a list of equations between formal composites of these symbols, which we intend as the relations that shall hold.
\end{enumerate}
Can we cook up a \SY{monoid} in which these---and no other---relations hold?
The answer is yes, and the recipe combines two constructions we have already seen: the free \SY{monoid}, and quotients by \SY{congruences} (\cref{sec:quotients-algebra}).

The idea is simple.
The free \SY{monoid}~$\listsof{\setA}$ contains all formal composites of the generator symbols, with no relations between them.
Our prescribed equations pick out certain pairs of lists that we want to merge, that is, a binary relation
\begin{equation}
    \relA \setsubseteq \listsof{\setA} \cartprod \listsof{\setA},
\end{equation}
where a pair~$\tup{\ela, \elb} \setin \relA$ encodes the equation ``$\ela = \elb$''.
Merging elements of a \SY{monoid} in a way that respects composition is exactly what quotients by \SY{congruences} do.
The only gap to bridge is that~$\relA$ itself is typically \emph{not} a \SY{congruence}: it records just the equations we explicitly wrote down, while a \SY{congruence} must also contain all their consequences (for instance, if~$\ela = \elb$ is imposed, then~$\elc \mtimes \ela = \elc \mtimes \elb$ must hold too, for every list~$\elc$).
The following lemma provides the bridge.

\begin{lemma}
    \label{lem:generated-congruence}
    Let~$\monoidA$ be a \SY{monoid} and let~$\relA \setsubseteq \monoidAset \cartprod \monoidAset$ be any binary endorelation on its underlying set.
    Among all \SY{congruences} on~$\monoidA$ that contain~$\relA$, there is a smallest one, namely the intersection of all of them.
    We call it the \maindef{congruence generated by}~$\relA$, and denote it by~$\equivalparam{\relA}$.
\end{lemma}

\begin{exercise}
    \label{ex:generated-congruence}
    Prove \cref{lem:generated-congruence}.
    That is, show that the intersection of any (non-empty) collection of \SY{congruences} on~$\monoidA$ is again a \SY{congruence}, and explain why at least one \SY{congruence} containing~$\relA$ always exists.
\end{exercise}
\begin{solution}
    Let~$\makeset{\equival_\setIel}_{\setIel \setin \setI}$ be a collection of \SY{congruences} on~$\monoidA$, and let~$\equival$ denote their intersection: $\monela \equival \monelb$ if and only if~$\monela \equival_\setIel \monelb$ for all~$\setIel \setin \setI$.
    Reflexivity, symmetry, and transitivity of~$\equival$ hold because they hold for each~$\equival_\setIel$ separately; the same argument applies to the compatibility condition \cref{eq:congruence-compatibility}.
    Finally, the collection of \SY{congruences} containing~$\relA$ is non-empty because the total relation (every pair of elements related, see \cref{ex:trivial-congruences}) is a \SY{congruence} and certainly contains~$\relA$.
\end{solution}

\begin{remark}
    There is also a useful ``bottom-up'' description of~$\equivalparam{\relA}$: two elements are related by~$\equivalparam{\relA}$ if and only if one can be transformed into the other by a finite sequence of steps, where each step replaces, inside a composite, a part matching one side of an imposed equation by the other side.
    This matches the intuitive picture of ``calculating with the relations'', as we did when simplifying~$\mapa^6 = \mapa^4$ at the beginning of this section.
\end{remark}

\begin{ctdefinition}[Presented monoid]
    \label{def:presented-monoid}
    Let~$\setA$ be a set (of ``generator symbols''), and let~$\relA \setsubseteq \listsof{\setA} \cartprod \listsof{\setA}$ be a binary relation (of ``imposed equations'').
    The \maindef{presented monoid} with generators~$\setA$ and relations~$\relA$, written~$\langle \setA \mid \relA \rangle$, is the quotient \SY{monoid}
    \begin{equation}
        \langle \setA \mid \relA \rangle \definedas \listsof{\setA} /_{\equivalparam{\relA}},
    \end{equation}
    where~$\equivalparam{\relA}$ is the \SY{congruence} generated by~$\relA$.
\end{ctdefinition}

\begin{remark}
    The analogous definition for \SY{semigroups} uses the free \SY{semigroup} of \emph{non-empty} lists (\cref{exa:string-semigroup}) in place of~$\listsof{\setA}$.
    For \SY{groups}, too, one can define presentations~$\langle \setA \mid \relA \rangle$; the construction requires a notion of \emph{free group}, in which formal inverses of the generators are also available as symbols.
    We will not pursue this here, but the underlying idea is the same.
\end{remark}

In examples, instead of writing out~$\relA$ as a set of pairs, we simply list the imposed equations.
We also write composites of generator symbols as strings (writing~$\setAel\setBel\setAel$ for the list~$\makelist{\setAel, \setBel, \setAel}$), and we write~$\emptylist$ for the empty list.

\begin{example}
    \label{exa:presentation-free-one-gen}
    The \SY{monoid}~$\langle \makeset{\setAel} \mid \; \rangle$, with one generator and no relations, is the free \SY{monoid} on one generator.
    Its elements are the classes of the lists
    \begin{equation}
        \emptylist, \quad \setAel, \quad \setAel\setAel, \quad \setAel\setAel\setAel, \quad \dots
    \end{equation}
    and since there are no relations, no two of these are merged.
    An element is thus determined exactly by the \emph{number} of occurrences of~$\setAel$, and composition adds these numbers.
    This is a description, by generators and relations, of the familiar \SY{monoid}~$\tup{\natnumbers, +, 0}$.
\end{example}

\begin{example}
    \label{exa:presentation-plant}
    Consider the \SY{semigroup} presented (as a semigroup) by
    \begin{equation}
        \langle \makeset{\setAel} \mid \setAel^5 = \setAel^4 \rangle,
    \end{equation}
    where~$\setAel^n$ abbreviates the list consisting of~$n$ copies of~$\setAel$.
    Using the imposed relation, any~$\setAel^n$ with~$n \geq 5$ can be shortened step by step to~$\setAel^4$, so the presented \SY{semigroup} has exactly the four elements~$\setAel$, $\setAel^2$, $\setAel^3$, $\setAel^4$ (and one can argue that no two of these are merged).
    This is an abstract description of the plant development \SY{semigroup} of \cref{exa:plant-trafo-semigroup}: the presentation captures the essential ``aging saturates'' behavior of the map~$\mapa$, while forgetting entirely what the states and the map ``are''.

    This example illustrates a general modeling point: a presentation records the \emph{interaction structure} of a system's actions---which sequences of actions are indistinguishable from which---without committing to any particular implementation of the states.
\end{example}

\begin{example}
    \label{exa:presentation-commutative}
    The \SY{monoid}
    \begin{equation}
        \langle \makeset{\setAel, \setBel} \mid \setAel\setBel = \setBel\setAel \rangle
    \end{equation}
    is the ``free \emph{commutative} monoid'' on two generators.
    Using the relation, the letters of any list may be sorted, moving all~$\setAel$'s to the front; so every element has a unique normal form~$\setAel^m \setBel^n$ with~$m, n \setin \natnumbers$.
    Composition adds the exponents componentwise.
    We already met this \SY{monoid} in \cref{exa:multiset-quotient}, in the guise of inventory bookkeeping: an element~$\setAel^m \setBel^n$ is a tally (``$m$ units of product~$\setAel$, $n$ units of product~$\setBel$''), and the imposed relation says exactly that the order of arrival is irrelevant.
\end{example}

\begin{example}
    \label{exa:presentation-toggle}
    The \SY{monoid}~$\langle \makeset{\setAel} \mid \setAel^2 = \emptylist \rangle$ has exactly two elements: $[\emptylist]$ and~$[\setAel]$.
    Indeed, the relation lets us delete pairs of adjacent~$\setAel$'s, so every list reduces to either the empty list or a single~$\setAel$, depending on the parity of its length.

    As an applied subexample, think of~$\setAel$ as pressing a toggle---the caps-lock key on a keyboard, or a simple on/off light switch.
    The relation~$\setAel^2 = \emptylist$ expresses that pressing the toggle twice is indistinguishable from not pressing it at all.
    Note that the resulting \SY{monoid} is actually a \SY{group}: $[\setAel]$ is its own inverse.
    Imposing relations can create inverses that were nowhere in sight in the free \SY{monoid}.
\end{example}

\begin{exercise}
    \label{ex:presentation-cyclic}
    Fix~$n \geq 1$.
    Show that the \SY{monoid}~$\langle \makeset{\setAel} \mid \setAel^n = \emptylist \rangle$ has exactly~$n$ elements, and that it is a \SY{group}.
    Which \SY{group} from \cref{sec:groups} do you recover for~$n = 4$?
\end{exercise}
\begin{solution}
    Using the relation, any power~$\setAel^m$ may be reduced to~$\setAel^{m \bmod n}$, so the elements are~$[\emptylist], [\setAel], \dots, [\setAel^{n-1}]$, and composition adds exponents modulo~$n$.
    Every element has an inverse: $[\setAel^k] \mtimes [\setAel^{n-k}] = [\setAel^n] = [\emptylist]$.
    For~$n = 4$ we recover the cyclic group of order 4 from \cref{exa:grp-Z4} (with~$[\setAel^k]$ corresponding to~$k$).
\end{solution}

\begin{remark}
    \label{rem:every-monoid-has-presentation}
    Every \SY{monoid} can be described by generators and relations, if all else fails in a rather uneconomical way: take \emph{all} of its elements as generators, and impose, for every pair of elements, the equation saying what their composite is (that is, transcribe the entire composition table into relations).
    The value of presentations lies at the other end of the spectrum: many important structures admit descriptions with very few generators and relations, and such descriptions are often the most compact and portable way to specify a structure---for a colleague, or for a computer.
\end{remark}

\begin{remark}
    \label{rem:word-problem}
    A note of caution.
    Given a presentation, deciding whether two lists of generators represent the \emph{same} element of the presented \SY{monoid} is known as the \emph{word problem}.
    In the examples above, the relations were simple enough that we could reduce every list to a normal form.
    In general, however, no algorithm can solve the word problem for all finite presentations---this is a celebrated undecidability result.
    So while presentations are a wonderfully compact specification format, extracting answers from them can be arbitrarily hard.
\end{remark}

